% ============================================================================
% Chapter 9 --- Compute and workload context
% ----------------------------------------------------------------------------
% Purpose : VMs, containers, orchestration, research clouds.
%           Belongs to the Execution plane (D5.1).
% Page target: ~5 pages.
% Dependencies: Ch. 5 (Kubernetes API, OpenStack API, OCI image, SLSA),
%               Ch. 7 (policy decisions enforced at admission and at runtime),
%               Ch. 8 (network attachment via CNI consistent with Ch. 8 zones),
%               Ch. 10 (storage volumes consumed by workloads).
% Mirrors layer chapter template established in Ch. 6 (D6.5).
% ----------------------------------------------------------------------------
% Decisions registered in _coherence-EN.md:
%   D9.1 minimum compliance set for the compute layer
%   D9.2 reference instantiation = OpenStack + Kubernetes + Harbor + SLSA
%   D9.3 alternatives = VMware + certified K8s, hybrid with cloud burst
%   D9.4 supply-chain provenance (SLSA) is part of the contract, not optional
%   D9.5 compute sovereignty scores
% ============================================================================

\chapter{Compute and workload context}
\label{ch:compute}

\begin{caveatbox}[Dependencies]
\noindent This chapter applies the layer template of
\trchap{ch:identity}{} to the compute layer. It consumes policy
decisions from \trchap{ch:policy}{} (admission and runtime), it
respects the network zone catalogue of \trchap{ch:network}{}
(\S\ref{sec:nw-reference}), and it produces workload-context events
consumed by \trchap{ch:observability}{}. The layer belongs to the
\emph{Execution plane}.
\end{caveatbox}

The compute layer is the place where institutional activities ---
research codes, teaching environments, administrative services, web
applications --- actually run. It is also a layer where lock-in
accumulates quietly: a custom hypervisor configuration, a
proprietary container registry, a cluster-management add-on whose
\gls{api} drifts from upstream Kubernetes can each, individually,
appear inexpensive, and together produce a stack from which exit
becomes a multi-year project.

%-----------------------------------------------------------------------------
\section{Functional role and interface contract}
\label{sec:co-contract}

The compute layer is responsible for: provisioning workloads as
virtual machines, containers, or composed workloads of both types;
enforcing admission decisions returned by the policy layer at the
moment of workload creation; running workloads under
network-attachment and storage-binding consistent with the activity
context; and producing the supply-chain provenance metadata that
allows the institution to verify what it actually runs.

\begin{contractbox}[Compute layer]
\noindent A compliant compute layer, in the sense of this
architecture, must:

\begin{itemize}[leftmargin=1.5em,nosep]
  \item expose a Kubernetes~\cite{kubernetes} API for container
    workloads, conformant to the upstream specification, without
    supplier extensions in critical paths;
  \item expose an OpenStack~\cite{openstack} API \emph{or} an
    equivalent standards-aligned IaaS API for virtual-machine
    workloads;
  \item accept and store container images in OCI image
    format~\cite{oci-image-spec}, in a registry whose access can be
    audited and exit-rehearsed;
  \item produce supply-chain provenance attestations conformant to
    the SLSA framework~\cite{slsa-spec} for every artefact produced
    internally;
  \item enforce admission decisions returned by \trchap{ch:policy}{}
    on every workload creation, with the decision recorded in the
    audit trail;
  \item integrate with the network layer
    (\trchap{ch:network}{}) through standard CNI plugins consistent
    with the institutional zone catalogue, and with the storage
    layer (\trchap{ch:storage}{}) through standard CSI plugins.
\end{itemize}
\end{contractbox}

The contract treats supply-chain provenance as a structural
requirement, not as a compliance afterthought. The reasoning is that
without verifiable provenance, the institution cannot make
auditable statements about what runs on its infrastructure --- a
gap that the data and institutional sovereignty dimensions surface
immediately.

%-----------------------------------------------------------------------------
\section{Reference instantiation: OpenStack + Kubernetes + Harbor + SLSA}
\label{sec:co-reference}

The reference instantiation combines OpenStack~\cite{openstack} for
IaaS, upstream Kubernetes~\cite{kubernetes} for container
orchestration, Harbor~\cite{harbor} as the OCI registry with image
scanning and signing (Cosign / Sigstore~\cite{sigstore}), and
SLSA-compliant~\cite{slsa-spec} build attestations emitted by the
institution's continuous-integration pipelines (typically built on
Tekton~\cite{tekton}).

\begin{instantiationbox}[Compute layer]
\noindent\textbf{Topology.} OpenStack provides multi-tenant
virtual-machine and bare-metal provisioning at the scale of the
institution; Kubernetes clusters --- one per major tenancy boundary
(administration, teaching, research projects of comparable scale) ---
run on top of OpenStack-provided infrastructure or on dedicated
metal where latency or isolation requires it. Container images live
in a Harbor registry deployed under the custody model the
institution has chosen for this layer, scanned at push time and
signed with cosign or an equivalent. Build pipelines (Tekton, Argo
Workflows, or comparable) emit SLSA-compliant provenance
attestations alongside the image. CNI plugins (Cilium is the
reference) implement the network policies derived from
\trchap{ch:network}{}; CSI plugins (Ceph-CSI is the reference)
expose the storage layer.
\end{instantiationbox}

\paragraph{Tenancy model.} Tenancy is the principal axis of
governance in this layer. The reference pattern uses Kubernetes
namespaces for fine-grained tenancy within a cluster, and separate
clusters for tenancy boundaries that warrant operational isolation
(cluster lifecycle independence, separate cryptographic identities,
distinct upgrade cadences). The choice between namespace tenancy
and cluster tenancy is a recurring design tension; this Reference
proposes that workloads of distinct sensitivity classifications
should not share a cluster.

\paragraph{Operational considerations.} The reference instantiation
is operationally non-trivial: OpenStack and Kubernetes each carry
substantial operational depth, and integrating them with the
identity, policy, network, and storage layers is the principal
ongoing engineering investment. An indicative envelope, for a
research-intensive institution running both an institutional
OpenStack and several Kubernetes clusters, would be the staffing
profile recorded for this layer in the institution's
infrastructure plan, scaled to the size and intensity of its
research compute portfolio. The labour market is broad
for Kubernetes itself and narrower for the integrated stack.

\begin{realworldbox}[Research-cloud precedents]
\noindent Indiana University's Jetstream2~\cite{jetstream2-2021},
operated under a national research-infrastructure funding programme,
provides a publicly documented research-cloud architecture built on
OpenStack and serving the multi-institutional research community.
Many institutions participate in a national or regional research
high-performance computing facility, several of which also rest on
OpenStack across multiple sites. Both should be
consulted as primary sources for design and operational lessons,
rather than as direct templates: each institution's compute layer
reflects its specific research mission and the context of its
national or regional research and education network, where one exists.
\end{realworldbox}

%-----------------------------------------------------------------------------
\section{Alternative~1 --- VMware with certified Kubernetes}
\label{sec:co-alt1}

\begin{alternativebox}[VMware-based compute]
\noindent The IaaS layer is provided by VMware vSphere~\cite{vmware-vsphere}; container
orchestration runs on VMware Tanzu (or another certified Kubernetes
distribution); image registry, build pipelines and CNI/CSI plugins
remain on the standard interfaces, with explicit care to avoid
captive extensions in their integration with the underlying
hypervisor.
\end{alternativebox}

\paragraph{Where it adds value.} A mature operational platform with
a deep skills market, broad enterprise documentation, and existing
integrations with most institutional tooling. For institutions whose
compute layer was built around vSphere over many years, the
operational momentum is substantial.

\paragraph{Where it constrains the institution.} VMware's
post-acquisition commercial trajectory has been the subject of
considerable public discussion, with reported changes to licensing
and pricing models that the institution would need to evaluate
against the financial-sovereignty dimension of
\trchap{ch:sovereignty-grid}{}. Beyond commercial considerations,
the technical constraint is the gradual coupling of the compute
layer to VMware-specific operators, controllers, and management
consoles whose absence would require non-trivial adaptation. The
architectural posture is to keep workloads expressed at the
Kubernetes API level and to treat VMware as the underlying
substrate, not as the API surface.

\paragraph{When it could be the right choice.} An institution with
a substantial existing investment in vSphere whose immediate
sovereignty priorities are not in the compute layer, and whose
financial planning has accounted for the supplier's pricing
trajectory. The choice would remain consistent with this
architecture provided that workloads are expressed in standard APIs
and that exit-rehearsal is practised periodically.

%-----------------------------------------------------------------------------
\section{Alternative~2 --- Hybrid with standards-based cloud burst}
\label{sec:co-alt2}

\begin{alternativebox}[Hybrid compute]
\noindent The institution operates a primary on-premises
open-source compute stack and federates additional capacity from
a public-cloud provider for bounded purposes (peak demand, specific
research workloads, geographically distributed collaboration). The
federation occurs through standard interfaces (Kubernetes API
across clusters, S3-compatible object stores) without proprietary
glue in the workload definition.
\end{alternativebox}

\paragraph{Where it adds value.} Elasticity for workloads whose
demand profile would over-provision an on-premises footprint, and
geographic flexibility for collaborations spanning multiple
jurisdictions. A hybrid posture also creates a continuously-exercised
exit rehearsal: workloads regularly cross the boundary between
on-premises and cloud, demonstrating that they are portable.

\paragraph{Where it constrains the institution.} Public-cloud
participation reintroduces the data-dimension considerations of
\trchap{ch:sovereignty-grid}{}: the data accessed by the cloud
portion of the workload falls under the cloud provider's regime.
The institutional governance of which workloads are eligible for
cloud burst, and under which conditions, becomes a permanent
operational concern. The hybrid posture is also more costly to
operate than a single-environment one.

\paragraph{When it could be the right choice.} An institution
with a research mission that includes large but episodic compute
demand (climate modelling, genomics, astronomical surveys, and
similar), where the alternative is either chronic
under-provisioning or institutional-scale over-provisioning, and
where the data-dimension governance is mature enough to support
per-workload eligibility decisions.

%-----------------------------------------------------------------------------
\section{Sovereignty grid applied}
\label{sec:co-sovereignty}

\begin{table}[H]
\centering\small
\caption{Per-dimension sovereignty profile of the compute
instantiations.}
\label{tab:co-sovereignty}
\begin{tabularx}{\textwidth}{%
  >{\hsize=0.6\hsize\linewidth=\hsize\bfseries\raggedright\arraybackslash}X%
  >{\hsize=1.2\hsize\linewidth=\hsize\centering\arraybackslash}X%
  >{\hsize=1.2\hsize\linewidth=\hsize\centering\arraybackslash}X%
  >{\hsize=1.2\hsize\linewidth=\hsize\centering\arraybackslash}X%
  >{\hsize=0.8\hsize\linewidth=\hsize\centering\arraybackslash}X%
}
\toprule
Dimension & Reference (OSS) & VMware-based & Hybrid (OSS + cloud) \\
\midrule
Data          & 3 Robust      & 3 Robust      & 2 Established \\
Technical     & 3 Robust      & 1--2 Partial  & 2 Established \\
Financial     & 2 Established & 1 Partial     & 2 Established \\
Operational   & 2 Established & 3 Robust      & 2 Established \\
Institutional & 3 Robust      & 1--2 Partial  & 2 Established \\
\bottomrule
\end{tabularx}
\end{table}

\paragraph{Driving indicators by dimension.} Data: storage location,
disclosure regime, especially for the cloud portion of the hybrid.
Technical: source visibility, presence of captive operators or
controllers in the cluster. Financial: cost predictability under
the supplier's announced pricing trajectory. Operational: skills
market depth and integration documentation. Institutional: roadmap
control over the hypervisor and the orchestrator.

\noindent The compute layer is the layer where the financial
dimension has been most visibly disturbed by recent supplier
consolidation in the proprietary segment. Institutions whose compute
strategy was set before that consolidation are reported to be
re-examining their position; this Reference treats that re-examination
as routine application of the sovereignty grid, not as an exceptional
event.

%-----------------------------------------------------------------------------
\section{Regulatory anchoring}
\label{sec:co-regulatory}

\noindent The compute layer carries the workload-context
attestations on which several security-of-processing obligations
of \loc{data-protection-regime} rest (integrity of processing,
posture verification of the host that processes personal data; in
the EU, \gls{gdpr}~Art.~32). It operationalises the
business-continuity and secure-development duties of
\loc{cybersecurity-baseline-regime} (in the EU, NIS2~Art.~21(2)(c),
through declarative redeployability, and (e), through SLSA-level
supply-chain attestations on the container images). Where the
institution adopts them, the layer also maps to ISO/IEC~27001~A.8.32
(change management) and NIST CSF \textsf{PROTECT} (PR.PS ---
platform security). Detailed mapping in
\trchap{ch:regulatory-mapping}{} \S\ref{sec:rm-gdpr},
\S\ref{sec:rm-nis2}, \S\ref{sec:rm-iso}, \S\ref{sec:rm-nist-csf}.

%-----------------------------------------------------------------------------
\section{Common pitfalls}
\label{sec:co-pitfalls}

\begin{pitfallbox}[Proprietary CNI or CSI in critical paths]
\noindent A Kubernetes deployment that uses a vendor-proprietary CNI
or CSI plugin for production workloads couples the cluster to that
vendor: an exit requires re-validation of every networked or
storage-bound workload. The mitigation is to use standards-aligned
plugins (Cilium for CNI, Ceph-CSI for CSI as references) for
production, and to confine vendor plugins to test or evaluation
clusters.
\end{pitfallbox}

\begin{pitfallbox}[Custom orchestrator above Kubernetes]
\noindent An institutional team writes a thin custom orchestrator on
top of Kubernetes to simplify a particular workflow. Within a few
years, the orchestrator has accumulated institutional knowledge that
the standard Kubernetes API does not capture, and migrating away
becomes expensive. The mitigation is to express workflow logic in
standards-aligned tools (Argo Workflows, Tekton) whose behaviour can
be re-implemented elsewhere if needed.
\end{pitfallbox}

\begin{pitfallbox}[Image registry lock-in]
\noindent Container images stored in a registry whose data export is
either undocumented or prohibitively slow create a structural exit
cost. The mitigation is to keep images in an OCI-compliant registry
whose export is exercised periodically as part of
\trchap{ch:conformance}{} (exit-rehearsal category).
\end{pitfallbox}

\begin{pitfallbox}[Hidden dependencies on supplier-specific operators]
\noindent Kubernetes admits operators --- in-cluster controllers
that extend the API --- contributed by suppliers, hardware vendors,
or open-source projects. Some of these operators carry implicit
dependencies on supplier services (telemetry endpoints,
licence-validation calls, configuration retrieval) whose absence
would silently degrade behaviour. The mitigation is to inventory
the operators in production clusters, classify their dependencies,
and either disable or substitute those whose dependencies erode the
sovereignty profile.
\end{pitfallbox}
